\chapter{Instruction Fine-Tuning and Parameter-Efficient Adaptation}\label{ch:instruct}

A pretrained model continues text; an assistant answers requests. If you ask base GPT-2
to ``Convert 45 kilometers to meters,'' it is just as likely to continue with a second
unit-conversion question as to answer. \emph{Instruction fine-tuning}, also called
supervised fine-tuning (SFT), closes that gap. We train the model with the ordinary
next-token objective on examples written as a prompt followed by the desired response,
so the most likely continuation of a prompt becomes its answer.

Algorithmically, SFT is pretraining with differently formatted data. The details are in
the data pipeline: how examples are templated, how variable-length sequences become
rectangular batches, which positions contribute to the loss, and how to pack many short
examples into one long row without letting them see each other. These details decide
whether SFT is correct and whether it is efficient, and they are the same machinery
behind the chat models served in \cref{ch:inference,ch:serving}. The second half of the
chapter turns to \emph{which} parameters to train. Low-rank adaptation (LoRA) and its
quantised variant QLoRA make fine-tuning a multi-billion-parameter model fit on one GPU.
They also make it possible to serve hundreds of task-specific variants from one base
model.

\begin{figure}[t]
\centering
\resizebox{\linewidth}{!}{%
\begin{tikzpicture}[
  stage/.style={draw=sysblue,fill=sysblue!8,rounded corners=2pt,minimum width=2.05cm,
    minimum height=1.05cm,align=center,font=\small},
  lstage/.style={stage,draw=napkin,fill=napkin!10},
  num/.style={circle,fill=sysblue,text=white,font=\scriptsize\bfseries,inner sep=1.5pt},
  arr/.style={-{Latex[length=2mm]},thick,sysblue!70}]
\node[stage] (s1) {Format\\instructions};
\node[stage,right=0.35cm of s1] (s2) {Batch:\\pad, shift,\\mask $-100$};
\node[stage,right=0.35cm of s2] (s3) {Data\\loaders,\\packing};
\node[stage,right=0.35cm of s3] (s4) {Load\\GPT-2\\medium};
\node[stage,right=0.35cm of s4] (s5) {Fine-tune};
\node[stage,right=0.35cm of s5] (s6) {Generate \&\\save\\responses};
\node[stage,right=0.35cm of s6] (s7) {Judge with\\a larger LLM};
\foreach \i in {1,...,7}{\node[num] at (s\i.north west) {\i};}
\foreach \a/\b in {1/2,2/3,3/4,4/5,5/6,6/7}{\draw[arr] (s\a) -- (s\b);}
\node[lstage,below=0.9cm of s3] (l1) {LoRA\\layers};
\node[lstage,right=0.35cm of l1] (l2) {QLoRA:\\NF4 base};
\node[lstage,right=0.35cm of l2] (l3) {Merge or\\multi-LoRA\\serving};
\node[num,fill=napkin] at (l1.north west) {8};
\node[num,fill=napkin] at (l2.north west) {9};
\node[num,fill=napkin] at (l3.north west) {10};
\draw[arr,napkin] (l1)--(l2); \draw[arr,napkin] (l2)--(l3);
\draw[densely dashed,napkin,thick] (l1.north) -- (s5.south west);
\node[font=\scriptsize\itshape,text=napkin,anchor=east] at ([xshift=-0.3cm]l1.west)
  {parameter-efficient track};
\end{tikzpicture}}
\caption{Chapter roadmap. Stages 1--7 build the instruction-tuning pipeline end to end.
Stages 8--10 replace full fine-tuning with low-rank adapters, quantise the frozen
base, and show how adapted models are deployed.}
\label{fig:instruct:roadmap}
\end{figure}

% =====================================================================================
\section{Introduction to instruction fine-tuning}\label{sec:instruct:intro}

Pretraining (\cref{ch:pretrain}) teaches a model the statistics of web text. Web text
contains questions followed by answers, but also questions followed by more questions,
forum signatures, and navigation menus. The pretrained model has no notion that a
request should be \emph{answered}. SFT shows it thousands of examples where it is.

\begin{figure}[ht]
\centering
\begin{tikzpicture}[
  stage/.style={draw=#1,fill=#1!12,rounded corners=2pt,minimum width=2.6cm,minimum height=0.85cm,
    align=center,font=\scriptsize},
  outbox/.style={draw=srcgray,fill=white,rounded corners=1pt,text width=4.1cm,font=\scriptsize\ttfamily,
    inner sep=4pt,align=left},
  arr/.style={-{Latex[length=2mm]},thick,srcgray}]
\node[stage=sysblue] (pre) at (0,0) {\textbf{Pretraining}\\next token on raw text};
\node[stage=keygreen] (sft) at (4.0,0) {\textbf{SFT}\\next token on\\prompt$\to$response};
\node[stage=labpurple] (pref) at (8.0,0) {\textbf{Preference / RL}\\reward on whole\\responses};
\draw[arr] (pre)--(sft); \draw[arr] (sft)--(pref);
\node[outbox] (o1) at (-0.5,-2.0) {Convert 45 kilometers to meters.\\Convert 12 miles to feet.\\Convert 3 liters to\dots};
\node[outbox] (o2) at (4.3,-2.0) {Convert 45 kilometers to meters.\\[2pt]\textbf{45 kilometers is 45,000 meters.}};
\draw[densely dashed,srcgray] (pre.south) -- (o1.north);
\draw[densely dashed,srcgray] (sft.south) -- (o2.north west -| sft.south);
\node[font=\scriptsize,text=srcgray,anchor=north] at (o1.south) {base model: continues the pattern};
\node[font=\scriptsize,text=srcgray,anchor=north] at (o2.south) {after SFT: answers the request};
\node[font=\scriptsize,text=labpurple,text width=2.6cm,align=center,anchor=north] at (8.6,-0.65)
  {\cref{ch:grpo,ch:grpo2,ch:rlsys}: optimise for \emph{which} answer is better};
\end{tikzpicture}
\caption{Where SFT sits. The same next-token loss is used throughout; what changes is the
data, and later the objective. Example outputs are illustrative.}
\label{fig:instruct:pipeline}
\end{figure}

Three facts frame the rest of the chapter.
\begin{enumerate}
\item \textbf{SFT needs surprisingly little data.} Strong results have come from tens of
thousands of examples, and in the LIMA study from about a thousand carefully curated
ones (Zhou et al., 2023). This is consistent with the ``superficial alignment''
hypothesis: the knowledge comes from pretraining, and SFT mostly teaches format and
behaviour. Our dataset has 1{,}100 examples.
\item \textbf{The model must learn to stop.} Every training response ends with an
end-of-text token. Without it, a fine-tuned model answers and then keeps generating
another instruction. This is why the collate function in \cref{sec:instruct:batches}
treats the first end-of-text token differently from the rest of the padding.
\item \textbf{Evaluation is the hard part.} Unlike classification, there is no single
correct string, so we will use a larger model as a judge (\cref{sec:instruct:eval}),
previewing the evaluation machinery of \cref{ch:eval}.
\end{enumerate}

% =====================================================================================
\section{Preparing a dataset for supervised instruction fine-tuning}\label{sec:instruct:data}

Our dataset is a JSON list of 1{,}100 records in the style popularised by
Alpaca (Taori et al., 2023). Each record has an \code{instruction}, an optional \code{input} that provides
context, and the reference \code{output}:

\begin{lstlisting}[style=plain]
{"instruction": "Rewrite the sentence in passive voice.",
 "input": "The committee approved the budget.",
 "output": "The budget was approved by the committee."}
{"instruction": "What is the boiling point of water at sea level in Celsius?",
 "input": "",
 "output": "The boiling point of water at sea level is 100 degrees Celsius."}
\end{lstlisting}

The tasks are short and varied: grammar fixes, unit conversions, definitions,
antonyms, simple factual questions. That makes them a good test of instruction
following on a 355M-parameter model that trains in minutes.

\subsection{Prompt templates}

Before tokenisation, every record is rendered into one string by a \emph{prompt
template}. The template defines the boundaries the model will learn: where the request
ends and the answer begins. \Cref{fig:instruct:styles} compares two common styles. The
\emph{Alpaca} style uses a fixed preamble and Markdown-like section headers. The
\emph{Phi-3} style uses compact role tags. Modern chat models use structured templates
of the second kind with dedicated special tokens (Qwen's ChatML,
\code{<|im\_start|>user \dots <|im\_end|>}, \cref{ch:inference}). Those tokens cannot
appear in user text by accident, which matters for both correctness and prompt-injection
resistance.

\begin{figure}[ht]
\centering
\begin{tikzpicture}[
  card/.style={draw=#1,fill=#1!6,rounded corners=2pt,text width=6.3cm,font=\scriptsize\ttfamily,
    inner sep=6pt,align=left},
  hl/.style={text=#1,font=\scriptsize\ttfamily\bfseries}]
\node[card=sysblue] (a) at (0,0) {%
Below is an instruction that describes a task. Write a response that appropriately completes
the request.\\[4pt]
\textcolor{sysblue}{\textbf{\#\#\# Instruction:}}\\Convert the temperature to Fahrenheit.\\[4pt]
\textcolor{sysblue}{\textbf{\#\#\# Input:}}\\25 degrees Celsius\\[4pt]
\textcolor{sysblue}{\textbf{\#\#\# Response:}}\\25 degrees Celsius is 77 degrees Fahrenheit.};
\node[card=keygreen,anchor=north west] (b) at ([xshift=0.5cm]a.north east) {%
\textcolor{keygreen}{\textbf{<|user|>}}\\Convert the temperature to Fahrenheit.\\25 degrees Celsius\\[4pt]
\textcolor{keygreen}{\textbf{<|assistant|>}}\\25 degrees Celsius is 77 degrees Fahrenheit.};
\node[font=\small\bfseries,anchor=south] at (a.north) {Alpaca style};
\node[font=\small\bfseries,anchor=south] at (b.north) {Phi-3 style};
\node[text width=6.3cm,font=\scriptsize,anchor=north west,text=srcgray] at ([yshift=-0.2cm]b.south west)
  {Fewer template tokens per example $\Rightarrow$ shorter sequences, faster training and
  inference, and more room for content in a fixed context window.};
\end{tikzpicture}
\caption{The same record rendered with two prompt templates. The response section is what
we ultimately want the model to produce; everything before it is the prompt.}
\label{fig:instruct:styles}
\end{figure}

\begin{lstlisting}[caption={Rendering a record with the Alpaca template (from \ident{llmsys/finetune.py}).},label={lst:instruct:alpaca}]
def format_alpaca(entry, with_response=False):
    text = ("Below is an instruction that describes a task. "
            "Write a response that appropriately completes the request."
            f"\n\n### Instruction:\n{entry['instruction']}")
    if entry.get("input"):                        # optional section
        text += f"\n\n### Input:\n{entry['input']}"
    text += "\n\n### Response:\n"
    if with_response:
        text += entry["output"]
    return text
\end{lstlisting}

The function returns the prompt alone (used at inference time) or prompt plus response
(used for training). Keeping one function for both prevents a subtle bug: a template
that differs by a single newline between training and inference measurably degrades a
fine-tuned model.

\begin{napkinmath}[What the template costs]
The Alpaca preamble and headers add roughly 30 GPT-2 tokens to every example. With
responses of about 20--60 tokens and instructions of about 10--30, the template is
roughly a third of all tokens processed. Training and evaluating with a compact
template cuts compute by about that fraction. In attention, where cost grows with
$T^2$, the saving is larger. The template also takes up context-window space that real
content could use.
\end{napkinmath}

\subsection{Splitting}

We shuffle with a fixed seed and split 85\% / 10\% / 5\% into training (935),
test (110) and validation (55) sets. The test set is larger than the validation set
because we will spend real effort on it (generation plus LLM judging). Validation is
used only to monitor loss.

% =====================================================================================
\section{Organizing data into training batches}\label{sec:instruct:batches}

In \cref{ch:classify}, the collate function only padded inputs. For SFT it has to do
more work. \Cref{fig:instruct:collate} shows the steps:

\begin{enumerate}
\item \textbf{Format and tokenise} each record (done once, in the dataset constructor).
\item \textbf{Append one end-of-text token} (id 50256 for GPT-2) to every sequence, so the
model learns to stop.
\item \textbf{Pad} every sequence in the batch to the length of the longest one plus one,
using the same id 50256.
\item \textbf{Shift}: inputs are positions $0..T-1$ of the padded row and targets are
positions $1..T$. This is exactly the shift used in pretraining.
\item \textbf{Mask padding in the targets}: every target position after the first
end-of-text token is replaced by $-100$, PyTorch's default \code{ignore\_index}, so the
loss ignores it. The \emph{first} end-of-text token stays a real target, because
predicting it is how the model learns to stop.
\end{enumerate}

\begin{figure}[ht]
\centering
\begin{tikzpicture}[
  tok/.style={draw=sysblue!60,fill=sysblue!12,minimum width=0.66cm,minimum height=0.5cm,
    font=\scriptsize\ttfamily,inner sep=0pt},
  eot/.style={tok,draw=keygreen,fill=keygreen!30},
  pad/.style={tok,draw=srcgray!60,fill=srcgray!15,text=srcgray},
  ign/.style={tok,draw=warnred!60,fill=warnred!15,text=warnred},
  ttl/.style={font=\scriptsize\bfseries,anchor=south west}]
\def\dx{0.68}
% group 1: padded rows
\node[ttl] at (-0.35,1.05) {steps 2--3: append EOT, pad};
\foreach \x/\t in {0/0,1/1,2/2,3/3,4/4}{\node[tok] at (\x*\dx,0.75) {\t};}
\node[eot] at (5*\dx,0.75) {E};
\foreach \x/\t in {0/5,1/6}{\node[tok] at (\x*\dx,0.2) {\t};}
\node[eot] at (2*\dx,0.2) {E};
\foreach \x in {3,4,5}{\node[pad] at (\x*\dx,0.2) {E};}
% group 2: inputs
\begin{scope}[xshift=5.4cm]
\node[ttl] at (-0.35,1.05) {step 4: inputs = row[:-1]};
\foreach \x/\t in {0/0,1/1,2/2,3/3,4/4}{\node[tok] at (\x*\dx,0.75) {\t};}
\foreach \x/\t in {0/5,1/6}{\node[tok] at (\x*\dx,0.2) {\t};}
\node[eot] at (2*\dx,0.2) {E};
\foreach \x in {3,4}{\node[pad] at (\x*\dx,0.2) {E};}
\end{scope}
% group 3: targets
\begin{scope}[xshift=10.2cm]
\node[ttl] at (-0.35,1.05) {steps 4--5: targets = row[1:], mask};
\foreach \x/\t in {0/1,1/2,2/3,3/4}{\node[tok] at (\x*\dx,0.75) {\t};}
\node[eot] at (4*\dx,0.75) {E};
\node[tok] at (0,0.2) {6};
\node[eot] at (\dx,0.2) {E};
\foreach \x in {2,3,4}{\node[ign] at (\x*\dx,0.2) {-100};}
\end{scope}
\draw[-{Latex[length=2mm]},thick,srcgray] (4.25,0.48) -- (4.95,0.48);
\draw[-{Latex[length=2mm]},thick,srcgray] (9.05,0.48) -- (9.75,0.48);
\node[font=\scriptsize,anchor=north west,align=left] at (-0.35,-0.25)
  {\textcolor{keygreen}{\rule{6pt}{6pt}} first EOT: a real target that teaches stopping\qquad
   \textcolor{srcgray}{\rule{6pt}{6pt}} padding (same id)\qquad
   \textcolor{warnred}{\rule{6pt}{6pt}} padding targets replaced by $-100$, ignored by the loss};
\end{tikzpicture}
\caption{The SFT collate function on two sequences \texttt{[0,1,2,3,4]} and
\texttt{[5,6]} (E = end-of-text id 50256). The pad token and the stop token share an id;
the target mask is what tells them apart.}
\label{fig:instruct:collate}
\end{figure}

\begin{lstlisting}[caption={The instruction collate function (from \ident{llmsys/finetune.py}).},label={lst:instruct:collate}]
IGNORE_INDEX = -100

def instruction_collate(batch, pad_id, ignore_index=IGNORE_INDEX,
                        allowed_max_length=None, mask_prompt=False, device="cpu"):
    items = [(b, 0) if not isinstance(b, tuple) else b for b in batch]
    max_len = max(len(ids) + 1 for ids, _ in items)       # +1 for the EOT
    inputs, targets = [], []
    for ids, prompt_len in items:
        row = list(ids) + [pad_id]                        # step 2
        row = row + [pad_id] * (max_len - len(row))       # step 3
        x = torch.tensor(row[:-1]); y = torch.tensor(row[1:])   # step 4
        y[len(ids):] = ignore_index                       # step 5: keep 1st EOT
        if mask_prompt and prompt_len > 0:
            y[:prompt_len - 1] = ignore_index             # optional: response-only
        if allowed_max_length is not None:
            x, y = x[:allowed_max_length], y[:allowed_max_length]
        inputs.append(x); targets.append(y)
    return torch.stack(inputs).to(device), torch.stack(targets).to(device)
\end{lstlisting}

\paragraph{Why $-100$ works.} With \code{reduction="mean"}, \code{F.cross\_entropy}
computes
\begin{equation}
\mathcal{L} = \frac{\sum_{i}\mathbb{1}[y_i \neq -100]\,\bigl(-\log p_\theta(y_i\mid x_{\le i})\bigr)}
                   {\sum_i \mathbb{1}[y_i\neq -100]},
\label{eq:instruct:masked}
\end{equation}
so ignored positions contribute neither to the numerator nor to the denominator. Adding
padding therefore leaves the loss value and its gradient exactly unchanged. That is the
property that makes dynamic padding safe. It also has a consequence that is easy to
miss: the mean is over \emph{tokens in the batch}, so long responses carry more total
weight than short ones, and the effective weight of an example depends on its
batch-mates. Gradient accumulation must reproduce this normalisation. Summing
per-microbatch means over-weights tokens in short microbatches; the fix is to sum
token losses and divide by the total target count of the full step.

\paragraph{Truncation.} \code{allowed\_max\_length} caps sequences at the model's
context window (1{,}024 for GPT-2). Truncate \emph{after} building targets so that the
input/target alignment survives. If truncation drops the response entirely, consider
filtering the example out instead, since it then contributes only prompt tokens.

\subsection{Masking the instruction}\label{sec:instruct:maskprompt}

By default the loss covers prompt tokens too, so the model is also trained to predict
the instruction text and the Alpaca preamble. A common variant masks every target that predicts a prompt token, so only the response
is learned (\cref{fig:instruct:maskprompt}). The trade-offs:
\begin{itemize}
\item \emph{Response-only loss} puts all gradient signal on what we want the model to
produce, and stops the model from learning to generate user-side text. Most modern
SFT recipes use it, especially for multi-turn chats, where only assistant turns are
trained.
\item \emph{Full-sequence loss} gives more training signal per example. Some studies
report it helps when responses are very short relative to prompts.
\end{itemize}
In the collate function, target index $j$ predicts token $j+1$, so masking the prompt
means masking targets $0..P-2$ for a prompt of $P$ tokens. Target $P-1$, the first
response token, is kept.

\begin{figure}[ht]
\centering
\begin{tikzpicture}[
  tok/.style={draw=sysblue!60,fill=sysblue!12,minimum width=0.85cm,minimum height=0.5cm,
    font=\scriptsize\ttfamily,inner sep=0pt},
  ign/.style={tok,draw=warnred!60,fill=warnred!15,text=warnred},
  rsp/.style={tok,draw=keygreen,fill=keygreen!25},
  lab/.style={font=\scriptsize,anchor=east}]
\node[lab] at (-0.3,0.6) {tokens};
\foreach \x/\t in {0/Below,1/is,2/\dots,3/\#\#\#,4/Resp.,5/:\textbackslash n}{\node[tok] at (\x*0.9,0.6) {\t};}
\foreach \x/\t in {6/45,7/km,8/is}{\node[rsp] at (\x*0.9,0.6) {\t};}
\node[lab] at (-0.3,-0.1) {targets};
\foreach \x in {0,...,4}{\node[ign] at (\x*0.9,-0.1) {-100};}
\foreach \x/\t in {5/45,6/km,7/is,8/\dots}{\node[rsp] at (\x*0.9,-0.1) {\t};}
\draw[decorate,decoration={brace,amplitude=4pt,mirror}] (-0.42,-0.45) -- (4.95,-0.45)
  node[midway,below=4pt,font=\scriptsize] {predicting prompt tokens: masked};
\draw[decorate,decoration={brace,amplitude=4pt,mirror}] (5.0,-0.45) -- (7.6,-0.45)
  node[midway,below=4pt,font=\scriptsize] {response: trained};
\end{tikzpicture}
\caption{Response-only training. The target that predicts the first response token
(here ``45'') sits at the last prompt position, which is why the mask stops at $P-1$.}
\label{fig:instruct:maskprompt}
\end{figure}

\begin{pitfall}[Tokenisation at the prompt boundary]
We compute $P$ by tokenising the prompt alone. BPE is not compositional: the tokens of
\texttt{prompt+response} are not always the tokens of \texttt{prompt} followed by those
of \texttt{response}, because a merge can span the boundary. The Alpaca template ends
in a newline, which no GPT-2 merge crosses, so it is safe. A template that ends in a
space, such as \texttt{"Answer: "}, is not. Production code tokenises the two segments
separately and concatenates the id lists, or uses a special-token boundary.
\end{pitfall}

\subsection{Packing: removing padding altogether}\label{sec:instruct:packing}

Bucketing (\cref{sec:classify:loaders}) reduces padding; \emph{packing} eliminates it.
We concatenate several EOT-terminated examples into one row of exactly $T$ tokens, and
fill rows with a bin-packing heuristic (first-fit decreasing is within $11/9$ of optimal
in the worst case and near-optimal in practice). Two corrections keep packing exactly
equivalent to training on the examples separately:

\begin{enumerate}
\item \textbf{Block-diagonal causal attention.} A token in document 2 must not attend to
document 1, so the attention mask is the causal mask AND a same-document mask.
\item \textbf{Per-document positions.} Position ids restart at 0 for each document.
Otherwise the second document is processed at positions it would never see at
inference time.
\end{enumerate}

\begin{figure}[ht]
\centering
\begin{tikzpicture}[x=0.36cm,y=0.36cm]
% padded batch
\node[font=\scriptsize\bfseries,anchor=south west] at (0,6.4) {padded: 4 rows};
\foreach \r/\n in {0/12,1/5,2/8,3/3}{
  \fill[sysblue!35] (0,5-\r*1.3) rectangle (\n,5.9-\r*1.3);
  \fill[srcgray!20] (\n,5-\r*1.3) rectangle (12,5.9-\r*1.3);}
\node[font=\scriptsize,anchor=west,text=srcgray] at (12.3,3.3) {44\% padding};
% packed
\node[font=\scriptsize\bfseries,anchor=south west] at (19,6.4) {packed: 3 rows};
\fill[sysblue!35] (19,5) rectangle (31,5.9);
\fill[keygreen!45] (19,3.7) rectangle (27,4.6); \fill[napkin!50] (27,3.7) rectangle (30,4.6);
\fill[labpurple!40] (19,2.4) rectangle (24,3.3); \fill[srcgray!20] (24,2.4) rectangle (31,3.3);
\node[font=\scriptsize,anchor=west,text=srcgray] at (31.3,3.3) {pad only in the tail};
% mask
\begin{scope}[shift={(19,-10.6)},x=0.36cm,y=0.36cm]
\node[font=\scriptsize\bfseries,anchor=south west] at (0,8.6) {mask for row 2 (doc A: 8, doc B: 3)};
\foreach \i in {0,...,10}{\foreach \j in {0,...,10}{
  \pgfmathtruncatemacro{\da}{\i<8 ? 1 : 2}
  \pgfmathtruncatemacro{\db}{\j<8 ? 1 : 2}
  \ifnum\j>\i \else\ifnum\da=\db
    \ifnum\da=1 \fill[keygreen!45] (\j*0.75,8-\i*0.75) rectangle ++(0.75,-0.75);
    \else \fill[napkin!50] (\j*0.75,8-\i*0.75) rectangle ++(0.75,-0.75); \fi
  \fi\fi
  \draw[white] (\j*0.75,8-\i*0.75) rectangle ++(0.75,-0.75);}}
\node[font=\scriptsize,anchor=west,text width=3.6cm] at (9,4)
  {blank = forbidden\\(future or other document)};
\end{scope}
\begin{scope}[shift={(0,-10.6)},x=0.36cm,y=0.36cm]
\node[font=\scriptsize\bfseries,anchor=south west] at (0,8.6) {position ids, row 2};
\foreach \p [count=\k from 0] in {0,1,2,3,4,5,6,7,0,1,2}{
  \pgfmathtruncatemacro{\d}{\k<8 ? 1 : 2}
  \ifnum\d=1 \fill[keygreen!45] (\k*1.2,6) rectangle ++(1.2,1); \else \fill[napkin!50] (\k*1.2,6) rectangle ++(1.2,1); \fi
  \node[font=\tiny] at (\k*1.2+0.6,6.5) {\p};}
\end{scope}
\end{tikzpicture}
\caption{Padding versus packing for four sequences of lengths 12, 5, 8 and 3 with
$T=12$. Packing needs a block-diagonal causal mask and position ids that restart per
document, so that documents are mathematically independent.}
\label{fig:instruct:packing}
\end{figure}

\begin{lstlisting}[caption={First-fit-decreasing packing and the block-diagonal mask.},label={lst:instruct:pack}]
def pack_sequences(seqs, max_len, pad_id, ignore_index=-100):
    rows, space = [], []
    for s in sorted(seqs, key=len, reverse=True):     # first-fit decreasing
        need = len(s) - 1                             # inputs s[:-1], targets s[1:]
        for r, free in enumerate(space):
            if free >= need:
                rows[r].append(s); space[r] -= need; break
        else:
            rows.append([s]); space.append(max_len - need)
    # ... fill input_ids, targets, segment_ids (1, 2, ...) and position_ids
    #     (restarting at 0) row by row; see llmsys/finetune.py

def segment_mask(segment_ids):                        # True = forbidden
    T = segment_ids.shape[1]
    causal = causal_mask(T, T, segment_ids.device)[None, None]
    other_doc = segment_ids[:, None, :, None] != segment_ids[:, None, None, :]
    return causal | other_doc
\end{lstlisting}

The lab verifies the equivalence numerically. The token-weighted mean loss of a packed
batch equals the loss obtained by running every example on its own, to within
$10^{-4}$.

\begin{systemslens}[Packing in real training stacks]
Materialising a $T\times T$ mask per row, as in \cref{lst:instruct:pack}, is fine for
teaching but wasteful at scale. Production kernels take \emph{cumulative sequence
lengths} instead: \code{cu\_seqlens} $=[0, \ell_1, \ell_1+\ell_2, \dots]$ for a flattened
batch with no padding at all. FlashAttention's \emph{varlen} kernels skip the
cross-document tiles entirely (\cref{ch:attention}), so packed attention costs
$\sum_i \ell_i^2$ instead of $T^2$ per row. For RoPE models the position ids are passed
alongside so each document's rotation restarts. Because of this:
\begin{itemize}
\item packing reduces MLP compute by the padding fraction, \emph{and} reduces attention
compute below even the unpadded-row cost $T^2$;
\item the per-step token count becomes constant, which stabilises memory use and step
time, and makes CUDA-graph capture possible (\cref{ch:inference});
\item the loss normalisation of \cref{eq:instruct:masked} becomes a per-token mean over
the whole pack. Schemes that weight each example equally need explicit per-example
weights.
\end{itemize}
Naive packing, which concatenates without the mask and position reset, is common in
pretraining, where documents are long and cross-document leakage is mostly harmless.
In SFT the examples are short and a model that sees previous answers can learn to copy
them, so do it properly.
\end{systemslens}

\begin{systemslens}[Where should the device transfer happen?]
A simple implementation moves each batch to the GPU inside the collate function, which is
convenient for single-process loading. With \code{num\_workers>0}, collate runs in
worker processes. Those should not touch CUDA: initialising CUDA in forked processes is
unsafe, and every worker would get its own CUDA context, costing hundreds of megabytes
each. The robust pattern is to keep collate on the CPU, set \code{pin\_memory=True} on
the \code{DataLoader}, and in the training loop call
\code{x.to(device, non\_blocking=True)}. Pinned host memory lets the copy run as an
asynchronous DMA that overlaps with the previous step's compute (\cref{ch:pytorch}).
\end{systemslens}

% =====================================================================================
\section{Creating data loaders for an instruction dataset}\label{sec:instruct:loaders}

The collate function has configuration (\code{pad\_id}, \code{allowed\_max\_length},
\code{mask\_prompt}) but \code{DataLoader} calls it with a single argument, the batch.
\code{functools.partial} binds the rest:

\begin{lstlisting}[caption={Instruction data loaders.},label={lst:instruct:loaders}]
from functools import partial
from llmsys.finetune import InstructionDataset, instruction_collate, format_alpaca

collate = partial(instruction_collate, pad_id=50256, allowed_max_length=1024,
                  mask_prompt=True)            # device stays "cpu": see Systems Lens
train_ds = InstructionDataset(train_data, tokenizer, fmt=format_alpaca)
train_loader = DataLoader(train_ds, batch_size=8, shuffle=True, drop_last=True,
                          collate_fn=collate, num_workers=2, pin_memory=True)
val_loader = DataLoader(InstructionDataset(val_data, tokenizer), batch_size=8,
                        shuffle=False, drop_last=False, collate_fn=collate)
\end{lstlisting}

\code{InstructionDataset} tokenises every example once in its constructor and stores
\code{(ids, prompt\_len)} pairs. Pre-tokenising pays off because tokenisation is pure
Python-side CPU work that would otherwise be repeated every epoch. For 1{,}100 examples
it takes a fraction of a second; for millions of examples, tokenisation becomes an
offline preprocessing job that writes binary shards (\cref{ch:text}).

Iterating over the training loader shows the defining property of dynamic padding:
batch widths vary (for example $8\times61$, $8\times76$, $8\times73$), each batch as
wide as its longest example.

% =====================================================================================
\section{Loading a pretrained LLM}\label{sec:instruct:load}

GPT-2 small (124M) is too weak to follow instructions well after a short SFT run. We
use \emph{GPT-2 medium} (355M parameters, 24 layers, $d=1024$, 16 heads) and load it
exactly as in \cref{ch:classify}, keeping the vocabulary head this time. Before
training, check what the base model does with a formatted prompt from the validation
set. Typically it repeats the \texttt{\#\#\# Response:} header, restates the
instruction, or starts a new instruction block. It recognises the format but has never
been trained to fill it in.

\begin{napkinmath}[Can I fully fine-tune GPT-2 medium on my GPU?]
GPT-2 medium has about $355$M unique parameters, so AdamW state is about
$16\times355\text{M}\approx\SI{5.7}{GB}$. Activations for a batch of 8 sequences of
about 80 tokens are small, a few hundred MB. A full fine-tune fits on any 8--12 GB GPU,
and even on a laptop CPU it runs in tens of minutes for two epochs over 935 examples.
At 7B parameters the same arithmetic gives \SI{112}{GB} for weights and optimizer state
alone. That is where \cref{sec:instruct:lora} begins.
\end{napkinmath}

% =====================================================================================
\section{Fine-tuning the LLM on instruction data}\label{sec:instruct:train}

The training loop is the one from \cref{ch:pretrain}. The only change is that the loss
uses \code{ignore\_index=-100}, which is \code{F.cross\_entropy}'s default, so the
pretraining \code{loss\_batch} works unchanged. We train for two
epochs with AdamW (learning rate $5\times10^{-5}$, weight decay $0.1$), evaluating
train and validation loss every few steps on a few batches.

\begin{lstlisting}[caption={SFT reuses the pretraining loop; only the data differs.},label={lst:instruct:train}]
from llmsys.train import train, loss_loader

model.to(device)
opt = torch.optim.AdamW(model.parameters(), lr=5e-5, weight_decay=0.1)
print("before:", loss_loader(model, train_loader, device, 5),
                 loss_loader(model, val_loader, device, 5))
hist = train(model, train_loader, val_loader, opt, device, num_epochs=2,
             eval_freq=5, eval_iter=5)
\end{lstlisting}

A typical run starts near a loss of 4 (the base model is surprised by the template) and
falls below 1 within the first epoch. Validation loss tracks training loss closely in
epoch 1 and starts to separate in epoch 2. Training longer gives diminishing returns
and eventually overfits a dataset this small. A useful habit is to print a generated
response for one fixed validation prompt after every epoch. A decreasing loss can hide
a model that has learned the template but not when to stop.

\begin{systemslens}[What dominates SFT wall-clock time]
For short SFT runs on small models the step itself is cheap, and time goes to
(i) per-step Python and kernel-launch overhead, (ii) evaluation passes, and
(iii) generation for qualitative checks, which runs at decode speed and is the slowest
thing in this pipeline. Once models grow, the levers of \cref{ch:pretrain} take over:
bf16 autocast, fused AdamW, activation checkpointing for long sequences, and
\code{torch.compile}. Packing becomes worth its complexity as soon as padding exceeds
roughly 20\%.
\end{systemslens}

% =====================================================================================
\section{Extracting and saving responses}\label{sec:instruct:extract}

To evaluate, we generate a response for every test instruction and store it next to the
reference. We format the prompt without the response, decode greedily with the KV-cached
generator from \cref{ch:inference} until the model emits the end-of-text token (or hits
a length cap), and strip the prompt.

\begin{lstlisting}[caption={Generating, extracting and saving test responses.},label={lst:instruct:extract}]
from llmsys.generate import generate
from llmsys.finetune import extract_response

def extract_response(generated_text, marker="### Response:"):
    tail = generated_text.split(marker)[-1]            # after the last header
    tail = tail.split("### Instruction:")[0]           # model started a new task
    return tail.replace("<|endoftext|>", "").strip()

results = []
for entry in test_data:
    ids = torch.tensor([tok.encode(format_alpaca(entry))], device=device)
    out = generate(model, ids, max_new_tokens=256, eos_id=50256)
    entry = dict(entry, model_response=extract_response(tok.decode(out[0].tolist())))
    results.append(entry)
json.dump(results, open("instruction-data-with-response.json", "w"), indent=2)
torch.save(model.state_dict(), "gpt2-medium355M-sft.pth")
\end{lstlisting}

Splitting on the response header is more robust than slicing off \code{len(prompt)}
characters: it survives tokeniser round-trip differences in whitespace. Cutting at a
new \code{\#\#\# Instruction:} header guards against models that have not fully
learned to stop.

\begin{systemslens}[This loop is a serving problem in miniature]
Generating 110 responses one at a time with a batch size of one leaves almost all of a
GPU idle, because each decode step reads every weight to produce one token
(\cref{ch:inference}). Batching the test prompts with left padding
(\code{generate\_batch}, \cref{ch:serving}) amortises each weight read over the whole
batch and finishes the same work many times faster. For thousands of prompts, use a
continuous-batching engine. The SFT and RL pipelines of later chapters
(\cref{ch:grpo,ch:distill}) are dominated by exactly this kind of bulk generation.
\end{systemslens}

% =====================================================================================
\section{Evaluating the fine-tuned LLM}\label{sec:instruct:eval}

There are three broad ways to evaluate an instruction-following model
(\cref{ch:eval} treats each in depth):
\begin{enumerate}
\item \textbf{Short-answer benchmarks} such as MMLU test knowledge with multiple-choice
questions. They are cheap and reproducible, but they barely test instruction following.
\item \textbf{Human preference comparisons}, such as the LMSYS Chatbot Arena, where
people vote between two anonymous models. They are the gold standard for
conversational quality, but slow and expensive.
\item \textbf{Automated conversational benchmarks} that use a strong LLM as the judge,
such as AlpacaEval and MT-Bench. They are scalable, but inherit the judge's biases.
\end{enumerate}

We follow the third route with a local judge. Ollama serves quantised open models
behind a small REST API. We run an 8B Llama~3 instruct model and ask it to score
every response from 0 to 100 against the reference.

\begin{figure}[ht]
\centering
\begin{tikzpicture}[box/.style={draw=#1,fill=#1!10,rounded corners=2pt,minimum height=0.9cm,
    text width=2.2cm,align=center,font=\scriptsize},arr/.style={-{Latex[length=2mm]},thick,srcgray}]
\node[box=srcgray] (t) at (0,0) {test record\\instruction, input,\\reference};
\node[box=sysblue] (g) at (3.0,0) {fine-tuned\\GPT-2 medium\\(generate)};
\node[box=keygreen] (e) at (6.0,0) {extract\\response};
\node[box=labpurple] (j) at (9.0,0) {judge LLM\\(Ollama, 8B)};
\node[box=napkin] (s) at (12.0,0) {parse score,\\aggregate};
\draw[arr] (t)--(g); \draw[arr] (g)--(e); \draw[arr] (e)--(j); \draw[arr] (j)--(s);
\draw[arr,densely dashed] (t.south) |- ++(0,-0.7) -| node[pos=0.25,below,font=\scriptsize]
  {reference answer goes into the judge prompt} (j.south);
\end{tikzpicture}
\caption{The evaluation pipeline. Generation is done once and saved; judging can then be
repeated with other judges or rubrics without touching the fine-tuned model.}
\label{fig:instruct:judge}
\end{figure}

\begin{lstlisting}[caption={A minimal Ollama client and a scoring prompt.},label={lst:instruct:ollama}]
def query_ollama(prompt, model="llama3", url="http://localhost:11434/api/chat",
                 seed=123, temperature=0.0, num_ctx=2048):
    payload = {"model": model, "messages": [{"role": "user", "content": prompt}],
               "options": {"seed": seed, "temperature": temperature,
                           "num_ctx": num_ctx}}
    req = urllib.request.Request(url, data=json.dumps(payload).encode(),
                                 headers={"Content-Type": "application/json"})
    chunks = []
    with urllib.request.urlopen(req) as resp:           # streamed JSON lines
        for line in resp:
            msg = json.loads(line)
            chunks.append(msg.get("message", {}).get("content", ""))
            if msg.get("done"):
                break
    return "".join(chunks)

def parse_score(text, lo=0, hi=100):
    m = re.search(r"-?\d+", text)
    return None if not m else max(lo, min(hi, int(m.group())))
\end{lstlisting}

The judge prompt (\code{build\_judge\_prompt} in \ident{llmsys/finetune.py}) shows the
instruction, the reference and the model answer, and asks for an integer only. We
average the parsed scores over the test set and count unparseable replies separately
rather than dropping them silently. A fine-tuned GPT-2 medium typically scores
around 50 on this scale. That is respectable for 355M parameters and far behind a
modern 8B instruct model.

\begin{pitfall}[LLM judges are measuring instruments, so calibrate them]
Judges prefer longer answers, answers in their own style, and the first of two options
in pairwise mode. Their scores drift between judge versions, and sampling makes them
non-deterministic unless temperature and seed are pinned (we pass
\code{temperature=0}, \code{seed=123}, and a fixed \code{num\_ctx}). Before trusting
a judge, score the reference answers themselves (they should be near 100) and a few
deliberately wrong answers (near 0), and report the judge model and version with every
number. \Cref{ch:eval} adds rubric-based grading and pairwise preference models.
\end{pitfall}

\begin{napkinmath}[What does judging cost?]
Each judgement is a prefill of about 200--300 tokens plus a few output tokens. An 8B
model in 4-bit on a laptop GPU processes prefill at thousands of tokens per second, so
110 judgements take well under a minute. Most of the wall-clock time is per-request
overhead and the sequential decode of the score. For thousands of evaluations, send
requests concurrently: Ollama and similar servers batch concurrent requests, which
raises throughput several-fold at the same per-request latency.
\end{napkinmath}

% =====================================================================================
\section{Conclusions and what comes next}\label{sec:instruct:conclusions}

At this point the from-scratch pipeline is complete: tokenisation and data sampling
(\cref{ch:text}), attention and the GPT architecture (\cref{ch:attention,ch:gpt}),
pretraining (\cref{ch:pretrain}), and two kinds of fine-tuning (this chapter and
\cref{ch:classify}). Modern assistants add a third training stage after SFT:
\emph{preference tuning}, with RLHF/PPO or direct preference optimisation (DPO), and
\emph{reinforcement learning with verifiable rewards} for reasoning. We build the
latter from scratch in \cref{ch:grpo,ch:grpo2} and the systems that run it at scale in
\cref{ch:rlsys}. Before that, Part~V turns the trained model into a service.

On staying current: the field moves quickly, but the core loop of data formatting,
masking, next-token loss and careful evaluation has stayed stable since the first
instruction-tuned models. New methods mostly change the data or the objective. The
best way to absorb them is the one this book follows: re-implement the core idea on a
small model, measure it, and only then reach for a framework.

% =====================================================================================
\section{Parameter-efficient fine-tuning with LoRA}\label{sec:instruct:lora}

Full fine-tuning updates every weight matrix $\mat{W}\in\R^{d_\text{out}\times d_\text{in}}$
by a dense $\Delta\mat{W}$. Low-rank adaptation (Hu et al., 2021) assumes the
\emph{update} has low intrinsic rank and parameterises it as a product,
\begin{equation}
\mat{W}' = \mat{W} + \Delta\mat{W} = \mat{W} + \frac{\alpha}{r}\,\mat{B}\mat{A},\qquad
\mat{A}\in\R^{r\times d_\text{in}},\ \mat{B}\in\R^{d_\text{out}\times r},\ r\ll\min(d_\text{in},d_\text{out}).
\label{eq:instruct:lora}
\end{equation}
$\mat{W}$ is frozen and only $\mat{A}$ and $\mat{B}$ are trained. Because of
distributivity we never form $\Delta\mat{W}$ during training:
\begin{equation}
\vect{y} = \vect{x}\mat{W}^\top + \frac{\alpha}{r}\,(\vect{x}\mat{A}^\top)\mat{B}^\top,
\end{equation}
which costs $r(d_\text{in}+d_\text{out})$ multiply-adds per token on top of the frozen
path. Three design choices matter:

\begin{description}
\item[Initialisation.] $\mat{A}$ is random (Kaiming-uniform) and $\mat{B}=\mat{0}$, so
$\Delta\mat{W}=\mat{0}$ at step 0. The adapted model starts \emph{exactly} at the
pretrained model, and the first gradient step reaches $\mat{B}$ through a non-zero
$\mat{A}$.
\item[Scaling $\alpha/r$.] The update is scaled by $\alpha/r$ so that changing $r$ does
not require re-tuning the learning rate. A common default is $\alpha = r$ or
$\alpha=2r$. Some implementations multiply by $\alpha$ directly, without
dividing by $r$, so their $\alpha=16, r=16$ corresponds to our $\alpha/r=16$. Watch for
this when porting hyperparameters between codebases. Our \code{LoRALinear} supports
both conventions through \code{scale\_by\_rank}.
\item[Placement.] The original paper adapted only the attention projections $W_q,W_v$.
Later work, notably QLoRA, found that adapting \emph{every} linear layer, MLPs
included, is needed to match full fine-tuning quality.
\end{description}

\begin{figure}[ht]
\centering
\begin{tikzpicture}[
  mat/.style={draw=#1,fill=#1!15,align=center,font=\small},
  arr/.style={-{Latex[length=2mm]},thick,srcgray},
  note/.style={font=\scriptsize,align=left}]
\node[mat=srcgray,minimum width=1.0cm,minimum height=0.5cm] (x) at (0,0) {$\vect{x}$};
\node[mat=sysblue,minimum width=2.5cm,minimum height=2.5cm] (W) at (-2.2,2.6)
  {$\mat{W}$\\[2pt]\scriptsize frozen\\\scriptsize $d_\text{out}\times d_\text{in}$};
\node[mat=keygreen,minimum width=2.5cm,minimum height=0.5cm] (A) at (2.2,1.6) {$\mat{A}$};
\node[mat=keygreen,minimum width=0.5cm,minimum height=1.7cm] (B) at (2.2,3.5) {$\mat{B}$};
\node[note,anchor=west] at (A.east) {$r\times d_\text{in}$, Kaiming init};
\node[note,anchor=west] at (B.east) {$d_\text{out}\times r$, zeros};
\node[circle,draw,inner sep=1pt,font=\small] (plus) at (0,5.3) {$+$};
\node[mat=srcgray,minimum width=1.0cm,minimum height=0.5cm] (y) at (0,6.2) {$\vect{y}$};
\draw[arr] (x.north) -- ++(0,0.45) -| (W.south);
\draw[arr] (x.north) -- ++(0,0.45) -| (A.south);
\draw[arr] (A.north) -- (B.south) node[midway,left,font=\scriptsize] {$r$};
\draw[arr] (W.north) |- (plus.west);
\draw[arr] (B.north) |- (plus.east) node[pos=0.8,above,font=\scriptsize] {$\times\,\alpha/r$};
\draw[arr] (plus) -- (y);
\node[note,text width=4.3cm,anchor=north west,text=srcgray] at (4.4,1.1)
  {trainable per layer: $r(d_\text{in}+d_\text{out})$.\\ E.g.\ $d_\text{in}=d_\text{out}=4096$, $r=16$:
   131K vs.\ 16.8M for $\mat{W}$ ($128\times$ fewer).};
\end{tikzpicture}
\caption{A LoRA-adapted linear layer. The frozen path and the low-rank path run in
parallel and are summed. After training, $\frac{\alpha}{r}\mat{B}\mat{A}$ can be folded
into $\mat{W}$ so inference costs nothing extra.}
\label{fig:instruct:lora}
\end{figure}

\begin{lstlisting}[caption={A LoRA linear layer with merge and unmerge (from \ident{llmsys/lora.py}).},label={lst:instruct:loralinear}]
class LoRALinear(nn.Module):
    def __init__(self, base: nn.Linear, rank=8, alpha=16.0, dropout=0.0,
                 scale_by_rank=True):
        super().__init__()
        self.base = base
        for p in self.base.parameters():
            p.requires_grad = False                     # frozen W (and bias)
        self.scaling = alpha / rank if scale_by_rank else alpha
        self.lora_A = nn.Parameter(torch.empty(rank, base.in_features))
        self.lora_B = nn.Parameter(torch.zeros(base.out_features, rank))
        nn.init.kaiming_uniform_(self.lora_A, a=math.sqrt(5))
        self.dropout = nn.Dropout(dropout) if dropout > 0 else nn.Identity()
        self.merged = False

    def delta_weight(self):
        return self.scaling * (self.lora_B @ self.lora_A)        # (d_out, d_in)

    def forward(self, x):
        y = self.base(x)
        if self.merged:
            return y
        low = F.linear(F.linear(self.dropout(x), self.lora_A), self.lora_B)
        return y + self.scaling * low

    @torch.no_grad()
    def merge(self):
        if not self.merged:
            self.base.weight += self.delta_weight(); self.merged = True

    @torch.no_grad()
    def unmerge(self):
        if self.merged:
            self.base.weight -= self.delta_weight(); self.merged = False
\end{lstlisting}

\subsection{Preparing the dataset}

We demonstrate LoRA on the spam classification task of \cref{ch:classify}
rather than on instruction data. That makes the comparison with full fine-tuning clean: the
same balanced SMS splits, the same dynamic-padding loaders, and the same accuracy
metric. Nothing about the data changes. The data pipeline is independent of how many
parameters are trained, so the loaders from \cref{lst:classify:collate} are reused
verbatim.

\subsection{Initializing the model}

We load pretrained GPT-2 small, swap in the two-class head
(\cref{lst:classify:head}), and check that the untrained classifier sits at chance
(about 46--49\% on the balanced splits). Then we freeze
\emph{every} parameter, the new head included. LoRA will supply all trainable
parameters. Counting trainable parameters before and after freezing (124M, then zero)
is a cheap guard against the most common bug, accidentally training the base.

\subsection{Parameter-efficient fine-tuning with LoRA}

\code{apply\_lora} walks the module tree and replaces every \code{nn.Linear} with a
\code{LoRALinear} wrapping it. A predicate can restrict which layers are adapted.

\begin{lstlisting}[caption={Replacing linear layers with LoRA layers.},label={lst:instruct:applylora}]
def apply_lora(model, rank=8, alpha=16.0, dropout=0.0, target=None,
               scale_by_rank=True, freeze_rest=True):
    if freeze_rest:
        for p in model.parameters():
            p.requires_grad = False
    def recurse(module, prefix):
        for name, child in list(module.named_children()):
            full = f"{prefix}.{name}" if prefix else name
            if isinstance(child, nn.Linear) and (target is None or target(full, child)):
                setattr(module, name, LoRALinear(child, rank, alpha, dropout,
                                                 scale_by_rank))
            elif not isinstance(child, LoRALinear):
                recurse(child, full)
    recurse(model, "")
    return model

model = add_classification_head(load_gpt2_small(), num_classes=2)
apply_lora(model, rank=16, alpha=16, scale_by_rank=False)   # LLM-book convention
print(count_trainable(model))      # (2666528, ...)
\end{lstlisting}

With $r=16$ on every linear layer of GPT-2 small plus the two-way head, the trainable
count is 2{,}666{,}528. Per block that is $4\times16\times(768+768)=98{,}304$ for the
attention projections and $16\times(768+3072)\times2=122{,}880$ for the MLP, times 12
blocks, plus $16\times(768+2)$ for the head. That is about 2.1\% of the model and
$47\times$ fewer than full fine-tuning. Training uses the same loop as
\cref{lst:classify:train}. With a learning rate around $5\times10^{-5}$ and five epochs,
LoRA reaches accuracy comparable to fine-tuning the last block, in somewhat more
wall-clock time. LoRA does not reduce compute much on a model this small (see the
next Systems Lens), and its benefits appear at scale.

\begin{keyidea}[Merging makes inference free]
After training, \code{merge()} folds $\frac{\alpha}{r}\mat{B}\mat{A}$ into $\mat{W}$,
and \code{merge\_lora} swaps every \code{LoRALinear} back to a plain \code{nn.Linear}.
The merged model has the original architecture and latency. The lab checks that
outputs are identical before and after merging, and that unmerging restores the base
weights. Keep adapters unmerged only when you need to switch between them quickly, as
in multi-LoRA serving below.
\end{keyidea}

\begin{systemslens}[What LoRA saves, and what it does not]
Write a frozen linear layer's training cost per token as forward $2P$, activation
gradient $\partial\vect{x}=\partial\vect{y}\,\mat{W}$ at $2P$, and weight gradient
$\partial\mat{W}=\partial\vect{y}^\top\vect{x}$ at $2P$, where $P=d_\text{in}d_\text{out}$.
LoRA removes only the last term: the activation gradient must still flow through
$\mat{W}$ to reach the adapters in lower layers. Training FLOPs per token therefore
drop from $\approx6N$ to $\approx4N$, at best $1.5\times$ cheaper. What LoRA mainly
removes is \emph{memory}:
\begin{itemize}
\item no gradients or Adam moments for the $N$ frozen parameters ($14N$ bytes saved);
\item frozen weights can be stored in bf16 without fp32 master copies, or quantised
(QLoRA, below);
\item checkpoints shrink from gigabytes to megabytes per task.
\end{itemize}
Activation memory is \emph{not} reduced. Every layer's input must still be kept for the
adapter weight gradients and for the backward pass through $\mat{W}$. Activation
checkpointing (\cref{ch:pretrain}) remains the tool for that.
\end{systemslens}

\begin{figure}[ht]
\centering
\begin{tikzpicture}
\begin{axis}[ybar stacked,width=0.82\linewidth,height=5.6cm,bar width=0.9cm,
  ylabel={GB (excluding activations)},symbolic x coords={Full FT,LoRA,QLoRA},xtick=data,
  ymin=0,ymax=120,legend style={font=\scriptsize,at={(0.98,0.97)},anchor=north east},
  nodes near coords={},enlarge x limits=0.25,grid=major,grid style={gray!20}]
\addplot[fill=sysblue!55,draw=sysblue] coordinates {(Full FT,13.48)(LoRA,13.48)(QLoRA,3.48)};
\addlegendentry{base weights (bf16 / NF4 + scales)}
\addplot[fill=warnred!50,draw=warnred] coordinates {(Full FT,94.36)(LoRA,0)(QLoRA,0)};
\addlegendentry{grads + fp32 master + Adam ($14N$)}
\addplot[fill=keygreen!55,draw=keygreen] coordinates {(Full FT,0)(LoRA,0.64)(QLoRA,0.64)};
\addlegendentry{adapters + their optimizer state}
\node[font=\scriptsize,anchor=south] at (axis cs:Full FT,107.8) {107.8};
\node[font=\scriptsize,anchor=south] at (axis cs:LoRA,14.1) {14.1};
\node[font=\scriptsize,anchor=south] at (axis cs:QLoRA,4.1) {4.1};
\end{axis}
\end{tikzpicture}
\caption{Persistent training memory for a 6.74B-parameter Llama-style model
($d=4096$, 32 layers, $d_\text{ff}=11008$), computed from the formulas in the text with
$r=16$ adapters on every linear layer (40.0M adapter parameters at 16 bytes each).}
\label{fig:instruct:membars}
\end{figure}

\subsection{QLoRA: a 4-bit frozen base}\label{sec:instruct:qlora}

Since the base is frozen, it only needs to be \emph{read}. QLoRA (Dettmers et al., 2023)
stores it in 4 bits and dequantises on the fly inside each matmul, while adapters and
activations stay in bf16. Three ingredients make this work without hurting quality:

\begin{enumerate}
\item \textbf{NormalFloat-4 (NF4).} Pretrained weights are roughly normally distributed,
so the 16 quantisation levels are placed at quantiles of $\mathcal{N}(0,1)$ rescaled to
$[-1,1]$, with an exact zero. With equal-probability bins every level is used about
equally often, which is the information-theoretically efficient choice for normal data.
Values are scaled per block of 64 weights by the block's absolute maximum.
\item \textbf{Double quantisation.} One fp32 scale per 64 weights adds $32/64=0.5$
bits per parameter. Quantising the scales themselves to 8 bits in blocks of 256
reduces this to $8/64 + 32/(64\cdot256)\approx0.127$ bits per parameter.
\item \textbf{Paged optimizers.} Adapter optimizer states live in unified memory, which
the driver can page to CPU RAM during memory spikes (long sequences) instead of failing
with an out-of-memory error.
\end{enumerate}

\begin{figure}[ht]
\centering
\begin{tikzpicture}
\begin{axis}[width=0.85\linewidth,height=4.6cm,xmin=-1.1,xmax=1.1,ymin=0,ymax=1.15,
  xlabel={normalised weight value $w/\text{absmax}$},ylabel={density (scaled)},
  ytick=\empty,grid=none,axis lines=left]
\addplot[thick,sysblue!60,domain=-1.1:1.1,samples=120,fill=sysblue!8] {exp(-(x*2.0)^2/2)} \closedcycle;
\foreach \v in {-1.0,-0.6962,-0.5251,-0.3949,-0.2844,-0.1848,-0.0911,0.0,0.0796,0.1609,0.2461,0.3379,0.4407,0.5626,0.7230,1.0}{
  \edef\tmp{\noexpand\draw[warnred,thick] (axis cs:\v,0) -- (axis cs:\v,{exp(-(\v*2.0)^2/2)});}\tmp
  \edef\tmp{\noexpand\fill[warnred] (axis cs:\v,0) circle (1.6pt);}\tmp}
\end{axis}
\end{tikzpicture}
\caption{The 16 NF4 code values (red) over a normal density, matching the shape of
block-normalised weights. Levels are dense near zero, where most weights live, and
sparse in the tails. Seven are negative and eight positive, so that zero is exact.}
\label{fig:instruct:nf4}
\end{figure}

\begin{lstlisting}[caption={Blockwise NF4 quantisation (one code per byte for clarity).},label={lst:instruct:nf4}]
def quantize_nf4(w, block=64):
    flat = w.detach().float().reshape(-1)
    pad = (-flat.numel()) % block
    if pad:
        flat = torch.cat([flat, flat.new_zeros(pad)])
    blocks = flat.view(-1, block)
    absmax = blocks.abs().amax(dim=1).clamp(min=1e-12)       # one scale / block
    normed = blocks / absmax[:, None]                         # in [-1, 1]
    codes = (normed[..., None] - NF4_LEVELS).abs().argmin(dim=-1).to(torch.uint8)
    return codes, absmax, tuple(w.shape)

def dequantize_nf4(codes, absmax, shape, dtype=torch.float32):
    vals = NF4_LEVELS[codes.long()] * absmax[:, None]
    return vals.reshape(-1)[:math.prod(shape)].reshape(shape).to(dtype)
\end{lstlisting}

\Cref{lst:instruct:nf4} stores one code per byte to keep indexing simple. Production
kernels (bitsandbytes) pack two codes per byte and fuse dequantisation into the GEMM,
so the 4-bit weights are expanded only in registers or shared memory. On random
Gaussian matrices our implementation reconstructs weights with a mean absolute error
of about 9\% of the mean magnitude, which the lab checks. \Cref{ch:efficient} develops
weight quantisation for inference, including GPTQ and AWQ, which choose levels using
calibration data rather than a fixed codebook.

\begin{napkinmath}[A 7B model on one 24 GB GPU]
Base weights in NF4 with double quantisation: $6.74\times10^{9}\times(4+0.127)/8
\approx\SI{3.5}{GB}$. Adapters ($r=16$, all linear layers): $40$M parameters
$\times$ 16 bytes $\approx\SI{0.64}{GB}$. That leaves roughly \SI{20}{GB} for
activations and the CUDA context, enough for batches of several thousand tokens with
activation checkpointing. Full fine-tuning of the same model needs about
\SI{108}{GB} before activations, more than one 80 GB GPU. QLoRA has one cost: every
forward and backward pass dequantises the weights, which typically makes each step
slower than LoRA on a bf16 base.
\end{napkinmath}

\subsection{Serving many adapters}\label{sec:instruct:multilora}

Merged adapters are free at inference, but a merged model serves exactly one task. A
platform hosting one base model and hundreds of customer-specific adapters faces a
different problem: requests for different adapters arrive interleaved, and one weight
copy per adapter does not fit. Systems such as Punica (Chen et al., 2023) and S-LoRA
(Sheng et al., 2023) keep adapters unmerged. They batch requests across adapters by
running the base GEMM once for the whole batch, then applying each row's adapter with a
gather-based batched kernel (Punica's ``BGMV''). S-LoRA additionally pages adapter
weights between host and GPU memory, in a unified pool with the KV cache.

\begin{figure}[ht]
\centering
\begin{tikzpicture}[
  row/.style={draw=srcgray,fill=srcgray!10,minimum width=1.9cm,minimum height=0.45cm,font=\scriptsize},
  arr/.style={-{Latex[length=1.6mm]},srcgray}]
\foreach \i/\a/\c in {0/2/keygreen,1/0/sysblue,2/none/srcgray,3/1/labpurple}{
  \node[row,fill=\c!15,draw=\c] (r\i) at (0,-\i*0.6) {req \i\ (adapter \a)};}
\node[draw=sysblue,fill=sysblue!12,minimum width=1.6cm,minimum height=2.3cm,align=center,font=\scriptsize]
  (W) at (3.4,-0.9) {shared\\$\vect{X}\mat{W}^\top$\\one GEMM};
\node[draw=napkin,fill=napkin!10,minimum width=2.4cm,minimum height=2.3cm,align=center,font=\scriptsize]
  (G) at (6.4,-0.9) {gather $\mat{A}_{a_i},\mat{B}_{a_i}$\\per row, then\\2 small bmm's};
\node[circle,draw,inner sep=1pt] (p) at (8.6,-0.9) {$+$};
\foreach \i in {0,...,3}{\draw[arr] (r\i.east) -- (W.west |- r\i.east);}
\draw[arr] (W) -- (G); \draw[arr] (G) -- (p);
\draw[arr] (W.north) |- ++(0,0.4) -| (p.north);
\node[font=\scriptsize,anchor=north] at (6.4,-2.25)
  {adapter stack: $\mat{A}\in\R^{K\times r\times d_\text{in}}$,\ $\mat{B}\in\R^{K\times d_\text{out}\times r}$};
\end{tikzpicture}
\caption{Batched multi-adapter inference. The expensive base matmul is shared by all
requests, and each row pays only for its own rank-$r$ correction. Row 2 uses the base
model only.}
\label{fig:instruct:multilora}
\end{figure}

\begin{lstlisting}[caption={A gather-based multi-LoRA layer (reference, not a fused kernel).},label={lst:instruct:multilora}]
class MultiLoRALinear(nn.Module):
    def __init__(self, base, n_adapters, rank, alpha=16.0):
        super().__init__()
        self.base, self.scaling = base, alpha / rank
        self.A = nn.Parameter(torch.zeros(n_adapters, rank, base.in_features))
        self.B = nn.Parameter(torch.zeros(n_adapters, base.out_features, rank))

    def forward(self, x, adapter_ids):          # x (B, T, d_in), ids (B,), -1 = none
        y = self.base(x)
        ids = adapter_ids.clamp(min=0)
        A, B = self.A[ids], self.B[ids]          # (B, r, d_in), (B, d_out, r)
        delta = torch.bmm(torch.bmm(x, A.transpose(1, 2)), B.transpose(1, 2))
        delta = delta * (adapter_ids >= 0).to(x.dtype)[:, None, None]
        return y + self.scaling * delta
\end{lstlisting}

\begin{napkinmath}[How many adapters fit?]
For a 7B Llama-style model with $r=16$ on the four attention projections only, one
adapter has $32\times4\times16\times(4096+4096)\approx16.8$M parameters, which is
\SI{34}{MB} in bf16. A thousand adapters are \SI{34}{GB}, too many to keep resident
next to the base weights and a KV cache on one 80 GB GPU. Hence S-LoRA's design: hot
adapters stay in GPU memory, cold ones live in host RAM and are fetched over PCIe on
demand. A 34 MB fetch at about \SI{25}{GB/s} takes about \SI{1.4}{ms}, cheap compared
with a prefill. The adapter math per token is
$2r(d_\text{in}+d_\text{out})$ FLOPs per adapted layer, about 0.25\% of the base model
here, so the overhead is dominated by gather and kernel efficiency rather than FLOPs.
\end{napkinmath}

\begin{deepdive}[The wider PEFT family]
LoRA dominates in practice, but the design space is larger.
\emph{Adapters} (Houlsby et al., 2019) insert small bottleneck MLPs between layers and
add inference latency unless they are fused. \emph{Prefix and prompt tuning} learn
virtual tokens and leave all weights untouched. They are cheap to store, but they cost
context length on every request. \emph{(IA)$^3$} learns per-channel rescaling vectors,
with even fewer parameters than LoRA. \emph{DoRA} (Liu et al., 2024) decomposes each
weight into a magnitude and a direction and applies LoRA to the direction, often
closing the remaining gap to full fine-tuning. The systems question is the same for all
of them: what must be stored per task, and can the per-task computation be batched
across tasks?
\end{deepdive}

% =====================================================================================
\section{Summary}\label{sec:instruct:summary}
\begin{itemize}
\item SFT is next-token training on templated prompt$\to$response data. The template
defines the model's interface and must be identical at training and inference time.
\item The collate function appends an end-of-text token, pads, shifts, and replaces
every padding target after the first EOT with $-100$. Optional prompt masking trains
only on responses. Masked positions drop out of both the numerator and the denominator
of the mean loss.
\item Packing removes padding entirely. To be exact it needs a block-diagonal causal
mask and per-document position ids; varlen attention kernels implement this without
materialising masks.
\item Keep collate on the CPU, pin memory, and transfer asynchronously. Bulk
generation for evaluation is a serving problem and should be batched.
\item LLM-as-a-judge evaluation scales, but the judge is an instrument that must be
pinned (model, version, temperature, seed) and sanity-checked.
\item LoRA trains $r(d_\text{in}+d_\text{out})$ parameters per layer, starts exactly at
the base model, and can be merged for free inference. It saves optimizer and gradient
memory, cuts training FLOPs by at most $1.5\times$, and does not reduce activation
memory.
\item QLoRA stores the frozen base in NF4 with double-quantised scales (about 4.13
bits per parameter), which puts 7B-model fine-tuning on one consumer GPU.
\item Unmerged adapters can be batched across tasks on a shared base (Punica, S-LoRA),
which enables multi-tenant serving.
\end{itemize}

% =====================================================================================
\section{Exercises}\label{sec:instruct:exercises}

\begin{exercise}[Phi-3 template]
Re-run SFT with \code{format\_phi3}. Compare the total number of training tokens, the
wall-clock time per epoch, and the judge score against the Alpaca template. Make sure
response extraction is adapted to the new template.
\end{exercise}

\begin{exercise}[Response-only loss]
Train with \code{mask\_prompt=True} and compare validation loss (careful: the losses are
averaged over different token sets, so compare on a common token set) and judge scores.
Under what data conditions would you expect full-sequence loss to win?
\end{exercise}

\begin{exercise}[Loss normalisation under gradient accumulation]
Show with a two-microbatch example that averaging per-microbatch mean losses differs
from the per-token mean over the whole step when microbatches have different numbers
of target tokens. Implement the correct normalisation in the training loop and verify
that the gradient matches a single large-batch step.
\end{exercise}

\begin{exercise}[Packing efficiency]
For the 935 training examples tokenised with GPT-2's BPE, compute padding waste for
(a) random batches of 8 with dynamic padding, (b) length-bucketed batches, and
(c) first-fit-decreasing packing into rows of 512 and 1{,}024 tokens. Then estimate the
attention FLOPs of (c) with a varlen kernel ($\sum\ell_i^2$) versus a dense mask
($T^2$ per row).
\end{exercise}

\begin{exercise}[Naive packing leaks]
Pack examples \emph{without} the block-diagonal mask and position reset. Construct a
probe showing that the model's prediction for a token in document 2 changes when
document 1 changes. How large is the effect after fine-tuning on naively packed data?
\end{exercise}

\begin{exercise}[Judge reliability]
Score the 110 reference answers themselves and 20 deliberately wrong answers with the
judge. Then re-score your model's responses with a second judge model, and report the
correlation between judges and the mean absolute score difference. What does this tell
you about the precision of a single reported score?
\end{exercise}

\begin{exercise}[LoRA rank sweep]
For LoRA on GPT-2 medium SFT, sweep $r\in\{1,4,16,64\}$ with $\alpha=2r$ and with fixed
$\alpha=16$. Plot judge score and validation loss against trainable parameters. Explain
why the $\alpha/r$ convention makes the learning rate more transferable across ranks.
\end{exercise}

\begin{exercise}[Where to put adapters]
Compare LoRA on (a) $W_q,W_v$ only, (b) all attention projections, (c) MLP only, and
(d) all linear layers at equal total trainable parameters (adjust $r$). Which wins on
the spam task, and which on SFT?
\end{exercise}

\begin{exercise}[FLOP accounting]
Instrument a LoRA training step with \code{torch.profiler} and attribute time to
forward GEMMs, activation-gradient GEMMs, weight-gradient GEMMs (adapters only) and
everything else. Compare with full fine-tuning and with the $6N$ versus $4N$
prediction.
\end{exercise}

\begin{exercise}[NF4 versus uniform int4]
Quantise each weight matrix of pretrained GPT-2 small with (a) NF4 and (b) uniform
symmetric int4, both with absmax scales per 64-weight block. Compare reconstruction
error and validation perplexity. Then check whether GPT-2's weights are actually close
to normal (plot a histogram of block-normalised weights).
\end{exercise}

\begin{exercise}[Multi-LoRA throughput]
Using \code{MultiLoRALinear}, measure the throughput of a batch of 64 rows that use
(a) one adapter, (b) 8 adapters, (c) 64 distinct adapters, against the base model alone.
Where does the gather cost show up, and how would a fused kernel avoid it?
\end{exercise}

% =====================================================================================
\section{Lab 17: The SFT data path, packing, LoRA and QLoRA}\label{sec:instruct:lab}

\begin{labbox}{17}{The SFT data path, packing, LoRA and QLoRA}
\labdir{lab17_instruct_lora}
\labsection{Goals}
Implement the parts of instruction tuning where correctness bugs hide (templating,
target construction, masking, packing) and the parameter-efficient machinery that
makes large-model fine-tuning feasible (LoRA, merging, batched multi-adapter layers,
NF4 quantisation). Verify each part against an exact reference, then train.

\labsection{Setup}
\code{cd code; LAB\_IMPL=starter python -m pytest -q labs/lab10\_instruct\_lora}.
Tests use a byte-level tokenizer and a tiny GPT (context 128). They run on CPU in under
a minute. \code{run.py} instruction-tunes GPT-2 medium (full or LoRA) on an Alpaca-style
JSON file, writes \code{responses.json}, and scores it with a local Ollama judge.

\labsection{Tasks}
\begin{enumerate}
\item \code{format\_alpaca}: exact template, optional input section.
\item \code{instruction\_collate}: append EOT, pad, shift, keep the first EOT, mask the
rest with $-100$, optional prompt masking and truncation. The test uses the
\texttt{[0..4]}, \texttt{[5,6]}, \texttt{[7,8,9]} example from
\cref{fig:instruct:collate}.
\item \code{pack\_sequences} and \code{segment\_mask}: first-fit-decreasing packing with
per-document positions and a block-diagonal causal mask. The test checks that the
packed loss equals the per-example token-weighted loss to $10^{-4}$.
\item \code{LoRALinear} (\code{delta\_weight}, \code{forward}, \code{merge},
\code{unmerge}), \code{apply\_lora}, \code{count\_trainable}: zero-init equivalence,
merge equivalence, exact trainable-parameter formula.
\item \code{MultiLoRALinear.forward}: gather-based batched adapters with a ``no adapter''
id. No Python loop over rows.
\item \code{quantize\_nf4} and \code{dequantize\_nf4}: blockwise absmax NF4 with padding
for sizes that are not multiples of the block.
\item \code{extract\_response} and \code{parse\_score}.
\item \code{finetune}: a short LoRA training loop. The test requires the loss to drop by
at least 30\% on a toy instruction task.
\end{enumerate}

\labsection{Measurements}
\begin{itemize}
\item Padding waste and attention FLOPs for padded, bucketed and packed batches on your
instruction data (Exercise 17.4).
\item On a GPU with \code{run.py}: peak memory, time per epoch and trainable parameters
for full fine-tuning versus LoRA ($r=16$) of GPT-2 medium. Expect LoRA to cut
persistent memory sharply and time per step only modestly. Compare judge scores.
\item Reconstruction error of NF4 versus int4 on real GPT-2 weights (Exercise 17.10).
\end{itemize}

\labsection{Deliverables}
Passing tests; a table of memory, time and score for full fine-tuning versus LoRA; one
paragraph explaining the measured memory and time differences using the $6N$/$4N$ and
$16N$/$2N$ models of this chapter.

\labsection{Stretch goals}
\begin{itemize}
\item Wrap the frozen base with \code{NF4LoRALinear} (provided in \ident{llmsys/lora.py})
and fine-tune GPT-2 medium QLoRA-style; compare quality and step time with bf16 LoRA.
\item Serve two LoRA adapters (for example spam classification and SFT) from one base
with \code{MultiLoRALinear} in every linear layer, and route a mixed batch.
\end{itemize}
\end{labbox}
